\documentclass[11pt]{article}
\usepackage[margin=1in]{geometry}
\usepackage{enumitem}
% =============================================================================
% Preambles/header.tex — shared LaTeX/Beamer preamble for this project
%
% Use from a Beamer lecture by prepending:
%     \input{header}
% and compiling with TEXINPUTS=../Preambles:$TEXINPUTS (the /compile-latex
% skill does this automatically).
%
% PALETTE CONTRACT. Color names defined here MUST match the names in
% Quarto/theme-template.scss so Beamer and Quarto renderings use the same
% palette. The scripts/check-palette-sync.sh script enforces this.
%
% To customize: edit the HEX values below AND the matching $variable in
% Quarto/theme-template.scss, then run scripts/check-palette-sync.sh.
% =============================================================================

% -----------------------------------------------------------------------------
% Engine detection + encoding
%
% Our /compile-latex skill uses XeLaTeX, where inputenc is unnecessary and
% can produce encoding warnings. Only load inputenc under pdfTeX.
% -----------------------------------------------------------------------------
\usepackage{iftex}
\ifPDFTeX
  \usepackage[utf8]{inputenc}
\fi

\usepackage{amsmath, amssymb, amsthm}
\usepackage{booktabs}
\usepackage{graphicx}

% -----------------------------------------------------------------------------
% PALETTE — mirrors Quarto/theme-template.scss variable names.
% Load xcolor and define colors BEFORE anything that references them
% (notably hyperref, hypersetup, and the Beamer color assignments).
% -----------------------------------------------------------------------------
\usepackage{xcolor}

\definecolor{primary-blue}{HTML}{012169}      % headings, accents
\definecolor{primary-gold}{HTML}{B9975B}      % emphasis, borders
\definecolor{highlight-yellow}{HTML}{F2A900}  % markers, alerts
\definecolor{light-bg}{HTML}{E8EDF5}          % title / section backgrounds
\definecolor{jet}{HTML}{1A1A1A}               % body text

% Semantic colors (used in diagrams and annotations)
\definecolor{positive}{HTML}{15803D}          % good / observed / identified
\definecolor{negative}{HTML}{B91C1C}          % bad / problematic / bias
\definecolor{neutral}{HTML}{525252}           % reference / context

% Highlight accents (match .hi-* classes in the SCSS)
\definecolor{hi-slate}{HTML}{314F4F}
\definecolor{hi-green}{HTML}{2E8B57}
\definecolor{hi-red}{HTML}{C41E3A}

% Semantic aliases so skill templates and snippets can write
% \textcolor{accent}{...} without hard-coding "primary-blue".
\colorlet{accent}{primary-blue}
\colorlet{accent2}{primary-gold}

% -----------------------------------------------------------------------------
% Hyperref — loaded AFTER xcolor + palette so link colors resolve.
% -----------------------------------------------------------------------------
\usepackage{hyperref}
\hypersetup{colorlinks=true, linkcolor=primary-blue, urlcolor=primary-blue,
            citecolor=primary-blue, pdfborder={0 0 0}}

% -----------------------------------------------------------------------------
% Course code — set once per deck (after \input{header}, before \begin{document})
% via \coursecode{CS401}. Shown in the footer on every slide so a deck's course
% is identifiable without opening the title page. Empty by default.
% -----------------------------------------------------------------------------
\makeatletter
\newcommand{\coursecode}[1]{\gdef\@coursecodetext{#1}}
\gdef\@coursecodetext{}
\makeatother

% -----------------------------------------------------------------------------
% Beamer theme assignments (only apply under Beamer).
%
% We detect Beamer via \@ifclassloaded{beamer}, which is the documented,
% non-internal way. This must be wrapped in \makeatletter/\makeatother
% because \@ifclassloaded contains an @-catcode character.
% -----------------------------------------------------------------------------
\makeatletter
\@ifclassloaded{beamer}{%
  \usefonttheme{professionalfonts}
  \setbeamercolor{structure}{fg=primary-blue}
  \setbeamercolor{frametitle}{fg=primary-blue}
  \setbeamercolor{title}{fg=primary-blue}
  \setbeamercolor{subtitle}{fg=primary-gold}
  \setbeamercolor{author}{fg=primary-blue}
  \setbeamercolor{institute}{fg=neutral}
  \setbeamercolor{date}{fg=primary-gold}
  \setbeamercolor{itemize item}{fg=primary-blue}
  \setbeamercolor{itemize subitem}{fg=primary-gold}
  \setbeamercolor{alerted text}{fg=primary-gold}
  \setbeamercolor{block title}{fg=white, bg=primary-blue}
  \setbeamercolor{block body}{bg=light-bg}
  % Semantic block variants: block=definition/concept (blue), exampleblock=
  % worked example (green/positive), alertblock=key takeaway or caution (gold).
  % Keep to <=2 colored boxes per slide across ALL three variants (INV-7).
  \setbeamercolor{block title example}{fg=white, bg=positive}
  \setbeamercolor{block body example}{bg=light-bg}
  \setbeamercolor{block title alerted}{fg=white, bg=primary-gold}
  \setbeamercolor{block body alerted}{bg=light-bg}

  % Minimal footer: course code (left) + page X / N (right), no navigation clutter
  \setbeamertemplate{navigation symbols}{}
  \setbeamertemplate{footline}{%
    \color{neutral}\footnotesize\hspace{0.7em}\@coursecodetext\hfill
    \insertframenumber/\inserttotalframenumber\hspace{0.7em}\vspace{0.3em}}
}{}
\makeatother

% -----------------------------------------------------------------------------
% TikZ setup — libraries the snippet gallery uses
% -----------------------------------------------------------------------------
\usepackage{tikz}
\usetikzlibrary{arrows.meta, positioning, calc, decorations.pathreplacing,
                fit, shapes.geometric, backgrounds}

% Shared TikZ styles — snippets and hand-written diagrams can pick these up
% via `\begin{tikzpicture}[every node/.style={...}]` or by naming specific
% styles (e.g., `\node[dag-node] (X) at (0,0) {X};`).
\tikzset{
  % Node styles for causal DAGs and similar
  dag-node/.style = {
    draw=primary-blue, fill=light-bg, rounded corners=2pt,
    minimum width=1.2cm, minimum height=0.9cm, align=center, font=\small
  },
  decision-node/.style = {
    draw=primary-gold, fill=white, diamond, aspect=1.8,
    minimum width=1.8cm, minimum height=1.2cm, align=center, font=\small,
    inner sep=1pt
  },
  % Edge styles — observed vs counterfactual
  observed-edge/.style = {-{Latex[length=2.5mm]}, thick, draw=jet},
  counterfactual-edge/.style = {-{Latex[length=2.5mm]}, thick, draw=neutral, dashed},
  confound-edge/.style = {-{Latex[length=2.5mm]}, thick, draw=negative, dashed},
  % Data-point styles for plots
  observed-dot/.style = {fill=primary-blue, draw=primary-blue, circle, inner sep=1.2pt},
  counterfactual-dot/.style = {fill=white, draw=neutral, circle, inner sep=1.2pt}
}

% -----------------------------------------------------------------------------
% Convenience macros
% -----------------------------------------------------------------------------
\newcommand{\muted}[1]{\textcolor{neutral}{#1}}
\newcommand{\key}[1]{\textcolor{primary-gold}{\textbf{#1}}}
\newcommand{\good}[1]{\textcolor{positive}{#1}}
\newcommand{\bad}[1]{\textcolor{negative}{#1}}

% Standout frame macro: full-bleed dark background for section transitions.
% Beamer-only (uses \begin{frame}/\setbeamercolor/\usebeamerfont) -- do not
% call from an article-class file (e.g. Notes/<CODE>/*-notes.tex). \muted,
% \key, \good, \bad, and \coursecode above are plain xcolor/gdef and are
% safe to use from either class; verified when Notes/article-class support
% was added.
% Expand: \transitionslide{Your transition title}
\newcommand{\transitionslide}[1]{%
  \begingroup
  \setbeamercolor{background canvas}{bg=primary-blue}
  \begin{frame}[plain]
    \centering
    \vfill
    {\usebeamerfont{title}\color{white}\Large #1\par}
    \vfill
  \end{frame}
  \endgroup
}

% Section divider frame (Beamer-only), an alternative to \transitionslide with a
% darker two-line style: near-black (jet) background, a small white label line
% above a larger gold title, and a thin gold rule along the bottom edge.
% Expand: \sectiondivider{Section 2}{Pointers}
\newcommand{\sectiondivider}[2]{%
  \begingroup
  \setbeamercolor{background canvas}{bg=jet}
  \begin{frame}[plain]
    \vspace*{3.0cm}
    {\color{white}\ttfamily\large #1\par}
    \vspace{0.5cm}
    {\color{primary-gold}\LARGE #2\par}
    \begin{tikzpicture}[remember picture, overlay]
      \draw[primary-gold, line width=2.5pt]
        ($(current page.south west)+(0,0.1)$) -- ($(current page.south east)+(0,0.1)$);
    \end{tikzpicture}
  \end{frame}
  \endgroup
}

% -----------------------------------------------------------------------------
% Channel watermark — "Concepts That Click" (YouTube).
%
% Article-class files (CompetitiveExam Q/A sets, Notes, Assignments) get a
% light diagonal draft watermark on every page plus a centered footer naming
% the channel. Beamer decks get a subtle TikZ background watermark instead
% (the footline already carries the course code). Centralized here so every
% MCQ PDF is branded without per-file edits.
% -----------------------------------------------------------------------------
\makeatletter
\@ifclassloaded{article}{%
  \usepackage{draftwatermark}
  \SetWatermarkText{Concepts That Click}
  \SetWatermarkScale{1}
  \SetWatermarkFontSize{2cm}
  \SetWatermarkLightness{0.86}
  \SetWatermarkAngle{45}
  \usepackage{fancyhdr}
  \pagestyle{fancy}
  \fancyhf{}
  \fancyfoot[C]{\footnotesize\textcolor{neutral}{Concepts That Click~|~YouTube: \href{https://www.youtube.com/@concepts.thatclick}{@concepts.thatclick}}}
  \renewcommand{\headrulewidth}{0pt}
  \renewcommand{\footrulewidth}{0pt}
  \fancypagestyle{plain}{%
    \fancyhf{}%
    \fancyfoot[C]{\footnotesize\textcolor{neutral}{Concepts That Click~|~YouTube: \href{https://www.youtube.com/@concepts.thatclick}{@concepts.thatclick}}}%
    \renewcommand{\headrulewidth}{0pt}%
    \renewcommand{\footrulewidth}{0pt}%
  }
}{}
\@ifclassloaded{beamer}{%
  \setbeamertemplate{background}{%
    \begin{tikzpicture}[remember picture, overlay]
      \node[rotate=45, opacity=0.055, align=center]
        at (current page.center)
        {\Huge\textcolor{neutral}{Concepts That Click}};
    \end{tikzpicture}%
  }
}{}
\makeatother 
\coursecode{JKSSB}
\renewcommand{\thesection}{58.\arabic{section}}
\newtheorem{example}{Example}[section]
\newenvironment{definitionbox}[1]{%
  \par\noindent\textbf{Definition (#1).}\ \itshape
}{\par\vspace{0.3em}}
\title{PYQ-Calibrated Final Retrieval and Audit --- Detailed Lecture Notes}
\author{JKSSB Computer Awareness}
\date{\today}

\begin{document}
\maketitle

\section{What this capstone is}
This lesson is the exam-eve mixed-retrieval capstone for the JKSSB Computer
Awareness course. It does not introduce new syllabus facts. It teaches the
candidate how to attack each JKSSB item format, how to audit negative stems,
how to handle defective options, and how to write an answer rationale that
names the trap in every distractor. The accompanying drill set holds the
course depth ratio: 50\% L1 direct recall, 35\% L2 contrast and
classification, 15\% L3 one-step application. All worked mini-examples below
are original items (\textit{[Original]}) built on concepts from the
Real-PYQ Evidence Registry and the trap registry; none reproduces an
official paper item verbatim.

\begin{definitionbox}{Retrieval at exam depth}
Answering a mixed set spanning all units under the real depth mix and the
real penalty (minus one-quarter mark per wrong answer, blank scoring zero),
using a fixed per-format attack and a written rationale for every item.
\end{definitionbox}

\begin{center}
\begin{tabular}{p{0.24\textwidth}p{0.66\textwidth}}
\toprule
\textbf{Term} & \textbf{Meaning in this lesson} \\
\midrule
Negative stem & A stem containing NOT, EXCEPT, incorrect, or CANNOT. \\
Near-neighbour & A distractor formed by swapping one word of a confusable pair. \\
Statement-evaluation & An item with statements I/II/III and combination options. \\
Defective option & Official option wording too loose to stand as fact. \\
Rationale & Why the key is right plus why each distractor is wrong. \\
\bottomrule
\end{tabular}
\end{center}

\section{Format attacks}
\subsection*{NOT/EXCEPT items}
Circle NOT, EXCEPT, or incorrect before reading the options. Mark each option
True or False in the margin and pick the odd one out. The three true options
are deliberately attractive, so never select the first true statement seen.
The worked pattern, after PYQ-03 (\textit{[Original]}): ``Which is NOT true
of cache?'' with options on recent data, speed, backup, and locality. The
backup option is the single false one, because cache holds copies of data
that live in RAM and is not a backup.

\subsection*{Near-neighbour items}
Retrieve the triple for every key term: what it is, what it is not, and the
nearest confusion. The productive pairs are RAM vs ROM, POP3 vs IMAP, Section
Break vs Page Break, digital signature vs encryption, MICR vs OCR vs OMR,
and compiler vs interpreter. The pattern after TRAP-15
(\textit{[Original]}): ``A digital signature encrypts the mail body so only
the recipient can read it'' is wrong, because a signature provides
authenticity and integrity while encryption provides confidentiality. One
swapped verb flips the item.

\subsection*{Abbreviation and shortcut items}
Say the full form on first sight: POP3 is Post Office Protocol version 3, ISP
is Internet Service Provider, MICR is Magnetic Ink Character Recognition. For
shortcuts, keep the exact key: F5 starts the slide show from the beginning,
F7 runs spell check, and an Excel formula must start with the equals sign.
Single-letter swaps (IMAP vs a coined variant, POP3 vs SMTP) are the standard
distractors, so spell each abbreviation out rather than recognising its shape.

\subsection*{Statement-evaluation items}
Judge each statement alone, writing T or F beside I, II, and III, and only
then keep the option whose combination matches. The pattern after PYQ-03 and
PYQ-04 (\textit{[Original]}): I. Registers are the fastest storage (true).
II. Cache is faster than main memory (true). III. RAM is non-volatile
(false). The surviving combination is I and II only. The false statement is
usually the volatility claim or the Program Counter claim.

\subsection*{Matching items}
Lock one certain pair first and strike out every option that breaks it. The
device drills decide most of these: the plotter and the sound card are
output, the scanner and the light pen are input, MICR reads magnetic ink, and
the BIOS initialises hardware and loads the operating system. The pattern
after PYQ-08 (\textit{[Original]}): among four device-to-role pairings, only
the MICR-to-magnetic-ink pairing survives. When two options survive, compare
only the pairs on which they differ.

\subsection*{Assertion--reason items}
Judge the assertion alone, judge the reason alone, and then ask whether the
reason explains the assertion. Both statements true but unlinked is the
standard trap. The pattern after PYQ-18 (\textit{[Original]}): the assertion
that IPv6 uses 128-bit addresses is true, and the reason that HTTP/3 runs
over QUIC rather than TCP is true, but the reason does not explain the
assertion, so the ``explains'' option must be rejected.

\begin{example}
A stem reads: ``Which of the following is NOT an input device?''
(\textit{[Original]}). Step 1: the format is NOT/EXCEPT, so the answer is
the single output device. Step 2: mark each option --- keyboard (input,
true), mouse (input, true), scanner (input, true), plotter (output, false).
Step 3: the plotter is the odd one out and the answer. Step 4: re-read the
stem to confirm the negative question was answered.
\end{example}

\section{The negative-stem audit method}
\begin{enumerate}
  \item Underline NOT, EXCEPT, incorrect, or CANNOT in the stem before
    reading any option.
  \item Convert the stem to its positive form: find the three true
    statements, because the leftover wins.
  \item Test every option with a written T/F mark; select the single False
    one.
  \item Re-read the stem once more before marking the OMR to confirm the
    negative question was answered.
\end{enumerate}
The discipline matters because the penalty is minus one-quarter mark per
wrong answer while a blank scores zero. A misread NOT-stem therefore costs
the mark plus the penalty. Guess only after eliminating at least two options.

\section{The defective-option audit method}
Some official options use wording too loose to stand as fact. The method,
per AP-02, is: quote the wording verbatim, flag the defect explicitly, and
teach the stable concept. Two instances recur in this course.

First, quoted verbatim: ``Program Counter stores the address of the
\textit{current} instruction being executed.'' Flag: the word \textit{current}
is defective. The stable concept is that the Program Counter holds the
address of the \textit{next} instruction to fetch (TRAP-04). Any option built
on the ``current instruction'' wording must be marked false.

Second, quoted verbatim from the option wording: ``BCC hides sender and
recipients from \textit{each other}.'' Flag: the phrase is defective. The
stable concept is that BCC hides recipients from \textit{each other}; the
sender always knows all recipients, and the To and Cc fields remain visible
to all (TRAP-16). Any option claiming the sender is hidden, or that BCC
works like encryption, must be marked false.

The Tier-1 final answer key always wins over any single option's wording, so
when wording and concept conflict, the concept above is what the candidate
carries into the hall.

\begin{example}
An option states ``ROM allows normal read and write''
(\textit{[Original]}, pattern after TRAP-02). Quote it, flag it: the words
``and write'' are defective. The stable concept is that ROM is read-only in
normal operation and writing to it needs a special process. The option is
false, and the rationale names the trap rather than memorising the sentence.
\end{example}

\section{Answer-rationale rules}
Every answer carries two halves: why the key is right, and why each
distractor is wrong. For each distractor, name the trap it instantiates (for
example, ``reverses RAM and ROM,'' ``swaps Section Break with Page Break,''
or ``states IPv6 as 256-bit''). Keep exam wording exact: abbreviations, full
forms, software names, protocol names, commands, and shortcut keys verbatim,
with explanatory prose in plain simple English. One fact per line. Never
invent a paper number, page, or key date. Never label a generated item as a
real PYQ; every item in the drill set is \textit{[Original]}.

\section{Per-unit quick-hit list}
Unit 1, Fundamentals and data representation: the smallest unit of data is
the bit; decimal 10 is binary 1010; ICs arrived in the third generation and
microprocessors in the fourth; BCD encodes each decimal digit in 4 bits; the
1024 ladder applies unless the option says otherwise.

Unit 2, CPU and architecture: the Program Counter holds the next
instruction's address; interrupts are handled at instruction boundaries, not
after the whole program; cycle counts vary by instruction; two's-complement
overflow is a wrong-sign result, not a carry out of the MSB; zero extension
is for unsigned widening; scenario items ask for the stated bottleneck, not
``slow CPU.''

Unit 3, Memory: RAM is volatile and ROM is non-volatile; cache is faster
than main memory and holds copies, not a backup; the CPU works through main
memory and never accesses secondary storage directly; the speed order is
registers, cache, RAM, secondary storage.

Unit 4, Input, output, and ports: the light pen, punch-card reader,
digitizer tablet, and scanner are input; the plotter, drum printer, and
sound card are output; MICR reads magnetic ink, OCR reads printed text, OMR
reads marks; large engineering drawings call for a plotter; the Section
Break, not the Page Break, enables different headers and footers.

Unit 5, Software and operating systems: a device driver is software and the
operating system is system software; the BIOS initialises hardware and loads
the operating system; open-source means the source can be viewed, changed,
and shared, which is distinct from freeware and shareware; multitasking runs
more than one program at a time.

Unit 6, Programming: machine, assembly, and high-level languages in order of
abstraction; assembler, compiler, and interpreter kept distinct; source code
vs object code; the stack follows last-in-first-out order.

Unit 7, Office and file formats: an Excel formula must start with the equals
sign; nested IF behaviour on blank and non-numeric input must be traced, not
guessed; COUNTIFS takes range-criterion pairs while SUMPRODUCT with
multiplication also counts under multiple conditions; Mail Merge combines
letters with an address database; a Select query retrieves data on a
condition; F5 starts the slide show from the beginning and F7 runs spell
check; the extension and the format are named separately.

Unit 8, Internet, web, and email: ISP is Internet Service Provider; the
at-sign separates the user name and the domain name; POP3 downloads mail to
the local machine while IMAP keeps it synced on the server; secure
transmission uses SSL/TLS; BCC hides recipients from each other.

Unit 9, Networking: IPv6 addresses are 128-bit with a fixed 40-byte header;
TCP orders delivery while UDP does not; ICMP is a control and diagnostic
protocol that transfers no application data; BGP routes between autonomous
systems (inter-domain) while OSPF and RIP work inside one domain; HTTP/3
uses QUIC over UDP, not TCP; the IaaS provider manages infrastructure while
the customer manages the operating system, runtime, applications, and data.

Unit 10, Cybersecurity: malware types kept distinct (virus, worm, Trojan,
ransomware, spyware, adware); a digital signature gives authenticity and
integrity while encryption gives confidentiality; phishing, spoofing, and
social engineering are recognition items; antivirus quarantine, firewall,
patching, backup, and multi-factor authentication each map to one control.

Unit 11, Governance and integration: e-NAM serves agricultural marketing;
G2C, G2G, G2B, and G2E expansions are direct recall; Digital India, NeGP,
e-Kranti, and NeGD roles are kept distinct; humanware, operator, and
administrator roles are named exactly.

\section{Exam retrieval and common confusions}
Ten distinctions prevent most errors on a mixed paper. RAM is volatile and
ROM is non-volatile. The Program Counter holds the next address, not the
current one. Cache is faster than main memory and is not a backup. Section
Break, not Page Break, changes headers and footers. A digital signature is
not encryption. POP3 downloads while IMAP syncs. IPv6 is 128-bit, never
256-bit. BGP is inter-domain. HTTP/3 uses QUIC over UDP. The IaaS provider
manages infrastructure, not the customer's operating system and data. Each
maps to a trap in the registry, and each appears in the drill set with its
distractors drawn from that trap.

\subsection*{Exercises}
\begin{enumerate}[label=\arabic*.]
  \item Convert the stem ``Which is NOT true of ROM?'' into its positive
    audit form, and state the selection rule.
  \item Quote verbatim one defective option discussed in this lesson, flag
    the defective word, and state the stable concept.
  \item For a statement-evaluation item with statements I, II, and III,
    describe the order in which the statements and the options are judged.
  \item For an assertion--reason item in which both statements are true but
    unlinked, state the correct verdict and explain why the ``explains''
    option is wrong.
  \item Write a two-half rationale for this item: ``IPv6 addresses are
    256-bit.'' State the verdict and name the trap.
\end{enumerate}

\subsection*{Solutions}
\begingroup\small
\begin{enumerate}[label=\arabic*.]
  \item Positive form: find the three true statements about ROM; the
    leftover false statement wins. Mark each option True or False and select
    the single False one, then re-read the stem before marking.
  \item Either ``Program Counter stores the address of the current
    instruction being executed'' (defective word: current; stable concept:
    the PC holds the next instruction's address) or ``BCC hides sender and
    recipients from each other'' (defective phrase: sender and recipients
    from each other; stable concept: BCC hides recipients from each other
    only).
  \item Judge each statement alone first, writing T or F beside I, II, and
    III; only then keep the option whose combination matches the T/F
    pattern.
  \item Verdict: both statements true, but the reason is not the explanation
    of the assertion. The ``explains'' option claims a direct link that does
    not hold, so it must be rejected even though each statement is true.
  \item Verdict: false. IPv6 addresses are 128-bit (TRAP-18); the option
    states the 256-bit distractor. The correction is one line: 128-bit
    addresses with a fixed 40-byte header.
\end{enumerate}
\endgroup
\end{document}
